\chapter{End-to-end data products}
\label{ch:end-to-end}

\begin{learningobjectives}
After this chapter, you should be able to:
\begin{itemize}
  \item define the endpoints and real boundaries of an end-to-end claim;
  \item trace one value through database, service, API, network, and client;
  \item distinguish build, deployment, serving, and monitoring responsibilities; and
  \item select tests and telemetry for client and server failure modes.
\end{itemize}
\end{learningobjectives}

\begin{chapterconnection}
This chapter is the Part I synthesis. It follows one value across every boundary
introduced earlier: typed representation, validation, reusable software,
tests, durable storage, an HTTP interface, a client, deployment, and operational
evidence. No single layer can protect the complete path by itself.
\end{chapterconnection}

\section{End to end is a scoped claim}

``End to end'' means from an explicitly named beginning to an explicitly named
end through the real boundaries claimed by a workflow or test. An instrument
event to a curated database row, a browser action to visible confirmation, and
a Git commit to a running artifact are different end-to-end paths.

An in-process API test with a temporary database is a valuable vertical-slice
test. It is not a deployed browser test through DNS, TLS, a load balancer,
application instances, and a managed database. Name the actual boundary rather
than using ``end to end'' as a synonym for comprehensive.

\begin{professionalbox}
An end-to-end notebook is a laboratory representation of a system, not the
operated system itself. It can make components and contracts visible, launch a
bounded local demonstration, and record observations. Production responsibility
belongs to independently testable modules, services, storage, deployment
workflows, and monitoring that continue to function after the notebook kernel
stops. The same separation strengthens research: a notebook can explain a
computational result while versioned data, configuration, code, and workflows
make the result independently reconstructible.
\end{professionalbox}

\section{One request across the system}

Consider an instrument submitting a temperature measurement:

\begin{center}
\begin{tikzpicture}[node distance=5mm]
  \node[flowbox] (client) {Instrument\\client};
  \node[flowbox,right=of client] (api) {HTTP API\\validation};
  \node[flowbox,right=of api] (service) {Domain\\service};
  \node[flowbox,right=of service] (repo) {Repository\\transaction};
  \node[flowbox,below=of repo] (db) {Database\\durable state};
  \node[flowbox,left=of db] (read) {Read API\\bounded query};
  \node[flowbox,left=of read] (browser) {Browser\\visible state};
  \draw[flowarrow] (client) -- (api);
  \draw[flowarrow] (api) -- (service);
  \draw[flowarrow] (service) -- (repo);
  \draw[flowarrow] (repo) -- (db);
  \draw[dataarrow] (db) -- (read);
  \draw[dataarrow] (read) -- (browser);
\end{tikzpicture}
\end{center}

The client sends JSON over HTTPS with identity, time, value, unit, and an
idempotency key. The API validates the public request model and HTTP semantics.
The domain service applies scientific or business policy. The repository owns
parameterized storage operations. The database enforces persistent constraints.
A later bounded query maps private rows into a public response that browser code
renders safely.

Each arrow is a contract. Ask about schema, units, identity, authentication,
authorization, retry behavior, ordering, latency, failure, and compatibility.

\section{Request anatomy}

For \code{GET /api/v1/measurements/latest?limit=20}, \code{GET} is the HTTP
method, \code{/api/v1/measurements/latest} is the path, and \code{limit=20} is a
query parameter. Method plus path identify an endpoint. Request headers carry
metadata. A request body can carry structured input. A response includes a
status code, headers, and optional result or error body.

The API is the agreement about these operations and semantics. JSON is one data
representation, HTTP is a protocol, and the server is a running implementation.

\section{Retries make writes ambiguous}

If a client times out after a POST request, it may not know whether the server
committed the write. Blind retry can create a duplicate. An idempotency key lets
the service recognize the same logical request and return the established
result. The key needs a defined scope, retention period, and transaction with
the effect it protects.

Explicit error models are part of the public API. Validation errors, conflicts,
authorization failures, unavailable dependencies, and internal defects should
produce stable machine-readable categories without exposing internal secrets.

\section{Frontend, backend, and hosting}

The \term{backend} implements server-side API, policy, orchestration, and
integrations. The \term{frontend} presents state and interaction, commonly in a
browser. A \term{client} may instead be a phone, instrument, script, or another
service.

A production request may pass through DNS, a content delivery network, web
application firewall, TLS termination, and a load balancer or API gateway before
reaching application compute. The application may use a managed database,
object store, queue, model service, and telemetry pipeline. Cloud products fill
roles; the architecture should be understandable without vendor names.

\section{CI/CD for production}

A release pipeline should build one immutable artifact, scan and test it,
deploy the same digest through environments, run migrations with compatible
application versions, and preserve rollback or roll-forward options.

Continuous delivery retains an explicit production decision. Continuous
deployment automates that decision for qualifying changes. Canary and blue-green
strategies reduce exposure; they do not replace compatibility design, database
recovery, or incident ownership.

\section{Monitor the client and server}

Server logs alone cannot explain a blank browser screen that failed before an
API request. A complete operational picture combines:

\begin{itemize}
  \item structured logs for discrete events and diagnostic context;
  \item metrics for rates, distributions, saturation, and alerting;
  \item traces for causal paths across services and dependencies;
  \item real-user monitoring for page, network, device, and JavaScript outcomes;
  \item synthetic journeys for controlled external checks; and
  \item privacy-reviewed retention and cardinality policies.
\end{itemize}

Correlation or trace identifiers connect evidence across boundaries. Do not use
personal data or secrets as identifiers. Service-level objectives translate
user-relevant reliability into measurable targets and response policy.

\section{Testing the real boundary}

Unit tests protect domain behavior. Repository integration tests exercise real
database semantics. Contract tests protect producer-consumer compatibility.
Vertical-slice tests send requests through API, service, and test database.
Browser component tests cover rendering and accessibility with controlled data.
Deployed journeys cover selected real network and hosting boundaries.

Phone and operating-system testing also has layers: responsive browser
viewports, browser engines, Android emulators, iOS simulators, cloud device
farms, and physical devices. Changing a user-agent string is not phone
emulation.

\begin{checkpointbox}
Define an end-to-end test for one measurement becoming visible on a phone. Name
the exact start, end, real dependencies, test data, cleanup, expected evidence,
and failures that remain outside the claim.
\end{checkpointbox}

\section{Worked example: trace one value without skipping a boundary}

An instrument records \code{21.7} degrees Celsius at 14:03:12 UTC. We will
follow that value until a scientist sees it in a browser.

\begin{workedexample}{measurement to visible evidence}
\textbf{1. Client representation.} The instrument constructs a request with
\code{experiment\_id}, a timestamp including its time-zone convention,
\code{value=21.7}, \code{unit="degC"}, and an idempotency key. Local validation
catches malformed fields but cannot authenticate the instrument to the server.

\textbf{2. Network request.} The client sends an authenticated HTTPS POST. TLS
protects transport to the configured endpoint; server-side authorization still
decides whether this identity may write to this experiment.

\textbf{3. Public API model.} The API parses JSON into typed fields, rejects
unknown or invalid input, and maps expected failures to stable HTTP responses.
It does not pass an arbitrary client dictionary directly to SQL.

\textbf{4. Domain service.} Application policy checks experiment state, allowed
units, plausible bounds, and duplicate-request behavior. Unit conversion, if
permitted, occurs in one named and tested place.

\textbf{5. Repository transaction.} The repository binds values to a
parameterized statement. The measurement and idempotency record commit in the
same transaction so a retry cannot produce a second logical observation.

\textbf{6. Read path.} A bounded query retrieves the stored record. A response
model exposes only public fields and states the unit explicitly.

\textbf{7. Browser rendering.} Client code checks the response status, decodes
JSON, escapes text, formats the timestamp, and renders \code{21.7 degrees C}.
Accessibility semantics ensure that the result is not communicated by color
alone.

\textbf{8. Operational evidence.} A correlation identifier connects safe log
events and traces across the request. Metrics count latency and failures. A
client error reporter can reveal a rendering defect that server telemetry
cannot see.
\end{workedexample}

At every stage, distinguish the value from its representation. The floating
point number in Python, bytes in JSON, a typed database column, and text in the
browser are not the same object. Contracts preserve identity, units, and
meaning across conversions.

\section{A timeout is an uncertainty about state}

Suppose the server commits the measurement, but the response is lost. The
client observes a timeout. It cannot infer that the transaction failed. If it
retries with a new request identity, the server may record a duplicate. If it
retries with the same idempotency key, the service can return the result of the
original logical request.

This is why retry policy cannot be added as a generic networking convenience.
Read operations are often safe to repeat, while writes require application
semantics. The database constraint, idempotency record, and effect must share a
transaction boundary or a crash can separate them.

\begin{misconceptionbox}
\textbf{Misconception:} end-to-end testing means testing everything, so one
large browser test can replace lower-level tests.

\textbf{Correction:} every end-to-end claim has named endpoints and omitted
risks. Broad tests give valuable integration evidence but are slower and less
diagnostic. Unit, integration, contract, vertical-slice, and deployed tests
protect different boundaries and should form a deliberate portfolio.
\end{misconceptionbox}

\section{From commit to production}

The following stages answer different questions:

\begin{longtable}{@{}p{0.19\textwidth}p{0.33\textwidth}p{0.38\textwidth}@{}}
\toprule
\textbf{Stage} & \textbf{Question} & \textbf{Typical evidence} \\
\midrule
\endhead
Build & Can source become an immutable artifact? & dependency lock, compiler or
package output, artifact digest \\
Verify & Does the candidate satisfy declared checks? & tests, type checks,
security and license scans \\
Deploy & Can that exact artifact enter an environment? & deployment record,
configuration validation, migration status \\
Release & Should users receive traffic now? & approval or policy decision,
canary health, feature flag state \\
Operate & Does the product meet user-relevant objectives? & metrics, traces,
logs, client outcomes, incident record \\
\bottomrule
\end{longtable}

Building twice for two environments weakens traceability because the binaries
may differ. Configuration can vary, but the verified artifact digest should
remain stable. Database changes require special care: old and new application
versions may overlap during rolling deployment, so schema evolution must be
compatible with both until the transition completes.

\section{Worked incident: the API is healthy but the page is blank}

Server dashboards show low latency and no increase in 5xx responses, yet users
report a blank results panel. This evidence narrows the hypothesis; it does not
prove the product is healthy.

Start with the user-visible path. Real-user monitoring shows a rise in browser
exceptions after frontend release \code{web-184}. Traces show that the API
returned 200 responses with a newly nullable field. A client log identifies a
formatting function that assumes the field is always numeric. The incident is a
consumer-contract failure: transport and backend serving succeeded while the
client could not render valid new data.

Recovery may disable the release or activate a feature flag. Prevention may
add a contract test for the nullable field, a component test for its visible
fallback, and a canary check that observes browser rendering. Each artifact
answers a specific question; ``add more monitoring'' is too vague to be an
action.

\conceptcheck{For the blank-page incident, what would logs, metrics, traces,
real-user monitoring, and a synthetic journey each contribute? Which signal is
closest to the user's actual outcome?}

\section{Exercises}

\subsection*{Foundational}
\begin{enumerate}
  \item Distinguish JSON, HTTP, an endpoint, an API, a server, and a client.
  \item Explain authentication versus authorization using the instrument
  request.
  \item Compare a log event, metric, and trace. Give one question each answers.
\end{enumerate}

\subsection*{Applied}
\begin{enumerate}
  \item Specify success and error responses for creating a measurement. Include
  validation, duplicate, authorization, and dependency failures.
  \item Design a vertical-slice test from HTTP request through a temporary real
  database. State what production boundaries it intentionally omits.
  \item Write a release checklist for a schema change that makes a response
  field nullable. Include producer, consumer, rollout, and rollback evidence.
\end{enumerate}

\subsection*{Design}
Define an end-to-end path for a scientific or industrial product of your
choice. Draw every process, network, data, and human boundary. For each arrow,
state the contract, likely failure, test layer, telemetry, authority, and
recovery owner. End with one precise service-level objective tied to a user
outcome.

\begin{chapterreview}
\textbf{Key ideas.} End to end is a scoped statement about real boundaries.
Meaning must survive representation changes. Retries can make write outcomes
ambiguous, so idempotency is a domain and transaction concern. CI/CD moves one
verified artifact through controlled environments. Operational evidence must
include the client-visible outcome.

\textbf{Vocabulary.} client; server; API; endpoint; HTTP method; status code;
JSON; authentication; authorization; idempotency; frontend; backend; artifact;
deployment; release; canary; log; metric; trace; real-user monitoring; SLO.

\textbf{Explain aloud.} How can every server dashboard look healthy while the
data product is unusable, and what evidence would reveal the gap?
\end{chapterreview}

\begin{notebooklink}
Lecture 8 notebook 00 implements a package-backed FastAPI, service, repository,
SQLite database, and responsive client as a bounded local laboratory. The
notebook explains and exercises the components; it is not the production
orchestrator. Companion files map the local vertical slice to containers,
CI/CD, browser testing, and monitoring artifacts.
\end{notebooklink}

\section{Part I takeaway}

A professional data scientist must be able to follow meaning across boundaries.
Part I began with a Python value and ends with an operated system. Models become
useful products only when data, interfaces, software, infrastructure, evidence,
and responsibility work together.
